\clearpage
\section{Wyznaczenia obciążenia powierzchni i obciążenia ciągu przyszłego samolotu}

    \subsection{Wykres ograniczeń}
        \paragraph{}
            Zgodnie z rodziałem 2.5 ksiąrzki \cite{galinski}. 
            Wykonałem wykres zależności odwrotności obciążenia ciągu od obciążenia powierzchni.

            \begin{figure}[!h]
                \centering \includegraphics[width=1\linewidth]{ObciazeniePowierzchni/ObciązeniePowierzchni.png}
                \caption{Wykres zależności czasu od roku produkcji}
            \end{figure}

        \paragraph{}
            Z wykresu jestem wstanie odzczytać iż dron by osiągnąć optymalny zasięg. 
            Musi cechować się natstępującytmi parametrami:
         
            \begin{align}
                \frac{W}{S} = 671
            \end{align}

            \begin{align}
                \frac{T}{W} = 0.319
            \end{align}

        \paragraph{} 
            Na podstawię odczytanych danych stwierdzam.
            Dron by osiądnąć optymalny zasięg potrzebuje powiechnie nośna powyżej 13m^2.
            Następnnie z wzoru 2.80 z \cite{galinski} obliczyłem,
            Układ napędowy powerplant powinnien cechowac się ciągiem powyżej 2757N.
            Silnik użyty do projektu powinien cechować sie mocą powyżej 145kW.